%! AP Statistics — Unit 5, Lesson 5.3 — Guided Notes (Answer Key)
\documentclass[10pt]{article}
\usepackage{apstats-article}
\usepackage{apstats-key}
\newcommand{\TallMath}[1]{$\displaystyle #1\rule[-1.4em]{0pt}{3.2em}$}

\begin{document}

\pageheader{Unit 5, Lesson 5.3}{Guided Notes --- Answer Key}
\namedateperiod

\vspace{0.1in}

\begin{objectivebox}
\textbf{Learning Targets}
\begin{enumerate}
  \item I can state the Central Limit Theorem and its conditions.
  \item I can compute $\mu_{\bar{x}}$ and $\sigma_{\bar{x}}$ for a sampling distribution of $\bar{x}$.
  \item I can determine when it is appropriate to use a normal model for $\bar{x}$.
\end{enumerate}
\end{objectivebox}

\begin{vocabbox}
\begin{description}
  \item[\termblanklong{Central Limit Theorem}] states that when a random sample of size $n$ is drawn from \emph{any} population with mean $\mu$ and standard deviation $\sigma$, the sampling distribution of $\bar{x}$ is approximately \ans{normal} when $n$ is sufficiently large.
  \item[\termblanklong{Standard error of $\bar{x}$}] the standard deviation of the sampling distribution: $\sigma_{\bar{x}} = \sigma/\sqrt{n}$.
  \item[\termblanklong{Independence condition}] each observation must be independent; when sampling without replacement, $n \le$ \ans{10}\% of the population.
\end{description}
\end{vocabbox}

\begin{hookbox}
\textbf{Hook — Candy Machine:} A machine fills bags with $\mu = 50$ candies and $\sigma = 12$.
\begin{itemize}
  \item Single bag: standard deviation = \ans{12}
  \item Mean of $n = 36$ bags: $\sigma_{\bar{x}} = 12/\sqrt{36} = $ \ans{2}
\end{itemize}
Why does averaging reduce variability?
\ansline{Averaging cancels out extremes:}\\
\ansline{large and small values offset, so $\bar{x}$ varies much less than single values.}
\end{hookbox}

\begin{notesbox}
\textbf{Central Limit Theorem — Three Key Facts}

\medskip
\begin{tabular}{@{} p{0.3\linewidth} p{0.6\linewidth} @{}}
\textbf{Center:} & $\mu_{\bar{x}} = $ \ans{$\mu$} \\[8pt]
\textbf{Spread:} & $\sigma_{\bar{x}} = $ \ans{$\sigma/\sqrt{n}$} \\[8pt]
\textbf{Shape:} & Approximately \ans{normal} when $n \ge 30$ OR when the population is \ans{normal} \\
\end{tabular}
\end{notesbox}

\begin{notesbox}
\textbf{Conditions for Using the Normal Model for $\bar{x}$} (UNC-3.H.3--H.5)

Check \emph{all three}:
\begin{enumerate}
  \item \textbf{Random:} the sample is a \ans{random sample (or randomized experiment)}.
  \item \textbf{Independence (10\% condition):} $n \le 10\%$ of \ans{the population}.
  \item \textbf{Normal/Large Sample:} population is normal, OR \ans{$n \ge 30$}.
\end{enumerate}

\medskip
\textit{Worked Example.} Household incomes: $\mu = \$58{,}000$, $\sigma = \$21{,}000$.
Sample of $n = 49$ households. Are conditions met?

\ansline{Random: assumed random sample of households.}\\
\ansline{10\%: 49 is less than 10\% of all households --- met.}\\
\ansline{Large sample: $n = 49 \ge 30$ --- met. All conditions are met.}

$\mu_{\bar{x}} = $ \ans{\$58{,}000} \hspace{1cm} $\sigma_{\bar{x}} = $ \ans{$21{,}000/\sqrt{49} = \$3{,}000$}
\end{notesbox}

\begin{practicebox}
\textbf{Quick Check}
\begin{enumerate}
  \item A population has $\mu = 80$, $\sigma = 15$, and is heavily left-skewed. For $n = 40$:
        \begin{enumerate}[label=\alph*.]
          \item $\mu_{\bar{x}} = $ \ans{80} \hspace{1cm} $\sigma_{\bar{x}} = $ \ans{$15/\sqrt{40} \approx 2.37$}
          \item Is the sampling distribution of $\bar{x}$ approximately normal? Why?
                \ansline{Yes: $n = 40 \ge 30$, so the CLT}\\
                \ansline{applies even though the population is left-skewed.}
        \end{enumerate}
  \item Same population, now $n = 9$. Is the sampling distribution of $\bar{x}$ approximately
        normal? Why or why not?
        \ansline{}\\
        \ansline{No: $n = 9 < 30$ and the population is left-skewed, so the CLT does not apply.}
\end{enumerate}
\end{practicebox}

\end{document}
